\documentclass{article}
\usepackage[T2A]{fontenc}
\usepackage[utf8]{inputenc}
\usepackage{amsthm}
\usepackage{amsmath}
\usepackage{amssymb}
\usepackage{amsfonts}
\usepackage[12pt]{extsizes}
\usepackage{fancyvrb}
\usepackage{indentfirst}
\usepackage[
  left=2cm, right=2cm, top=2cm, bottom=2cm, headsep=0.2cm, footskip=0.6cm, bindingoffset=0cm
]{geometry}
\usepackage[english,russian]{babel}

\begin{document}
\section*{Вариант 22}

Рассмотрим дискретную временную шкалу \(T = \{t_0 < t_1 < \dots < t_N\}\), в узлах которой заданы значения двух показателей: \(\{q(t_k)\}\) (например, объёма спроса, объёма ВВП и пр), и \(\{p(t_k)\}\) (цены, объёма трудовых или капитальных ресурсов и пр.), \(k = 0, \dots, N\). Считаем, что \(\displaystyle \min_{k \in 0,\dots,N} p(t_k) = p(t_0) < \min_{k \in 1,\dots,N} p(t_k)\) или \(\displaystyle \max_{k \in 0,\dots,N} p(t_k) = p(t_0) > \max_{k \in 1,\dots,N} p(t_k)\), \(\displaystyle \min_{k \in 0,\dots,N} q(t_k) = q(t_0) < \min_{k \in 1,\dots,N} q(t_k)\) или \(\displaystyle \max_{k \in 0,\dots,N} q(t_k) = q(t_0) > \max_{k \in 1,\dots,N} q(t_k)\).

Введём в рассмотрение нормированные индексы, соответственно, для \(k = 0, \dots, N\), взяв в качестве базы индексирования начальные значения показателя в точке \(t_0\):

\[
I_q(t_k) = q(t_k)/q(t_0),
\]

\[
I_p(t_k) = p(t_k)/p(t_0).
\]

Вычислим точечную эластичность показателя \(q\) относительно показателя \(p\) для каждого \(k = 0, \dots, N\) по формуле:

\[
E(t_k) = \frac{I_q(t_k) - 1}{I_p(t_k) - 1} = \frac{\Delta q/q}{\Delta p/p}.
\]

\end{document}